% =============================================================
\documentclass[11pt]{article}
\usepackage[a4paper,margin=2cm]{geometry}
\usepackage[spanish]{babel}
\usepackage[utf8]{inputenc}
\usepackage{amsmath,amssymb}
\usepackage{booktabs}
\pagestyle{empty}
\setlength{\parindent}{0pt}
\setlength{\parskip}{4pt}

% ---------- Datos del taller: editar aquí ----------
\newcommand{\curso}{Física Computacional (106018C)}
\newcommand{\taller}{Taller 5 --- Método de Monte Carlo y convergencia}
\newcommand{\semana}{Semana 5}
\newcommand{\lenguaje}{Fortran 2008 / Python}
\newcommand{\estudiante}{Cristian Andrés Cárdenas Muñoz (2518531)}
\newcommand{\fechaentrega}{26/10/2026} % Revisar fecha real de entrega
% ----------------------------------------------------

\begin{document}
\begin{center}
{\large\bfseries \curso} \\[2pt]
{\bfseries \taller} \\[2pt]
\semana\ \quad --- \quad Lenguaje: \lenguaje \\[2pt]
\estudiante\ \quad --- \quad \fechaentrega
\end{center}
\vspace{4pt}\hrule\vspace{6pt}

\textbf{1. Contexto físico.}
En física computacional, el muestreo aleatorio permite aproximar integrales y promedios de sistemas donde evaluarlos directamente puede resultar costoso. Como caso de referencia con solución conocida, se estimó el área del círculo unitario (\(\pi\)) y la integral de \(\sin x\) en \([0,\pi/2]\), para examinar precisión y convergencia.

\textbf{2. Método.}
Se usó Monte Carlo de acierto y error para el área del círculo y los métodos de valor medio (VM) y acierto y error (AE) para la integral. Se hicieron 20 corridas por cada \(N=10^2,\ldots,10^7\) en el cálculo de \(\pi\), y 1000 corridas de 1000 muestras para cada estimador de la integral:
\begin{equation*}
\hat\pi=4\frac{N_{\rm dentro}}{N},\qquad
\hat I_{\rm VM}=\frac{\pi}{2N}\sum_{i=1}^N\sin x_i,\qquad
s\propto N^{-1/2}.
\end{equation*}
La pendiente de \(\log_{10}s\) frente a \(\log_{10}N\) se ajustó por mínimos cuadrados en Fortran. La misma estimación de \(\pi\) se repitió con el generador LCG \((57,1,256)\) para examinar su sesgo.

\textbf{3. Resultado.}
Las cifras son medias de las corridas, excepto el LCG, que usa una sola secuencia periódica.
\begin{center}
\begin{tabular}{lccc}
\toprule
Caso & Valor obtenido & Referencia & Error relativo \\
\midrule
\(\pi\), \(N=10^7\) & 3.141530 & 3.141593 & 0.00198\% \\
Integral (VM) & 0.999731 & 1 & 0.02695\% \\
Integral (AE) & 1.000085 & 1 & 0.00852\% \\
\(\pi\), LCG \(N=10^7\) & 3.093750 & 3.141593 & 1.52288\% \\
\bottomrule
\end{tabular}
\end{center}
La pendiente fue \(m=-0.5079\) (referencia: \(-0.5\)). Aunque la media de AE quedó algo más cerca de 1, su dispersión fue mayor: \(s_{\rm AE}=0.024645\) frente a \(s_{\rm VM}=0.015292\).

\textbf{Verificación (PASS/FAIL):} 5 de 5 controles pasaron: convergencia, pendiente, los dos estimadores de la integral y el generador LCG.

\textbf{4. Obstáculo o decisión de implementación.}
Una media cercana al valor exacto no demuestra que el algoritmo sea estable. Por ello se repitieron los cálculos para medir la dispersión y se verificó la pendiente por un procedimiento independiente. También se acumularon conteos y sumas sin guardar millones de coordenadas, reduciendo el uso de memoria; fijar las semillas permitió reproducir la revisión.

\textbf{5. Conclusión.}
Monte Carlo se aproxima a los valores de referencia, pero mejorar la precisión estadística exige muchas más muestras. El valor medio fue más estable para esta integral; además, el LCG mostró que aumentar \(N\) no corrige un error sistemático por correlaciones. La calidad del generador importa tanto como la cantidad de puntos.

\textbf{6. Uso de herramientas de IA.}
Se utilizó ChatGPT para ordenar y direccionar ideas, corregir dudas puntuales en codigo, revisar una redaccion neutra y formal. Las propuestas se revisan y ajustan personalmente antes de la entrega.

\vspace{4pt}\hrule\vspace{4pt}
{\small Código fuente y notebook: repositorio de GitHub, carpeta \texttt{Taller\_5}.}
\end{document}
